% profiles/university-of-north-carolina-at-chapel-hill.tex — UNC-Chapel Hill
% ============================================================================
% source: https://gradweb.unc.edu/content/academics/thesis-diss/guide/format.html
% source: https://gradweb.unc.edu/content/academics/thesis-diss/guide/accessibility.html
%         (The Graduate School, "Thesis and Dissertation Guide", sections II and
%         III. The Guide states that where it disagrees with a style manual,
%         "the regulations set forth in this Guide take precedence".)
% accessed: 2026-09-02
% checked:  top-20 research-institution pass
%
% VERIFIED and set here: geometry, chapter-opening margin, spacing, and the
% accessibility requirements -- which are the most specific in this whole set.
% LEFT TO THE NEUTRAL FALLBACK: title page and approval page (the Guide gives
% them as sample pages, which were not read line by line on this pass).
%
% A TRAP FOR THE NEXT READER: the Guide's own "Print complete guide" link points
% at pdf/ETDGuide2013.pdf -- a 2013 document that predates the entire
% accessibility section. Use the live web pages, not that PDF.

\thesisinstitution{University of North Carolina at Chapel Hill}

% "All copies of a thesis or dissertation must have the following uniform
% margins throughout the entire document: Left: 1" (or 1 1/4" to ensure
% sufficient room for binding the work if desired); Right: 1"; Bottom: 1" (with
% allowances for page numbers); Top: 1"."
\thesissetgeometry{left=1in, right=1in, top=1in, bottom=1in}

% The exception, quoted in full because it is the reason \thesissetchapteropening
% exists: "The first page of each chapter (including the introduction, if any)
% begins 2" from the top of the page. Also, the headings on the title page,
% abstract, first page of the dedication/acknowledgments/preface (if any), and
% first page of the table of contents begin 2" from the top of the page."
%
% This covers the chapter openings and the \chapter*-based front-matter pages.
% It does NOT reach the title page, which UNC specifies through a sample page.
\thesissetchapteropening{2in}

% "The text must appear in a single column on each page and be double-spaced
% throughout the document." Exceptions: "Blocked quotations, notes, captions,
% legends, and long headings must be single-spaced throughout the document and
% double-spaced between items."
\thesissetspacing{double}

% --- Accessibility -----------------------------------------------------------
% UNC is the only school in this set that names a WCAG version and level, and it
% states the obligation as legal rather than stylistic.
\thesisrequiretagging
\thesisrequiredocumentlanguage
\thesisaccessibilityrequirement{Policy}
	{"All students submitting theses or dissertations must ensure their work
	 meets digital accessibility standards. This requirement is not only a
	 matter of university policy, it's also a legal obligation under the updated
	 ADA Title II web" accessibility rule.}
\thesisaccessibilityrequirement{WCAG 2.2 AA}
	{UNC has adopted WCAG 2.2 Level AA "for all digital materials", and the
	 Guide names it as the standard a thesis is measured against.}
\thesisaccessibilityrequirement{Headings}
	{Real heading structure is required. "Do not use styled text (e.g., bold or
	 large font) as a substitute for proper headings."}
\thesisaccessibilityrequirement{Alt text}
	{Required for all images and non-text content. Avoid "image of..." -- screen
	 readers announce that already. For complex images such as charts, give
	 brief alt text plus a long description near the image or in an appendix
	 referenced from the alt text.}
\thesisaccessibilityrequirement{Tables}
	{Add alt text or place descriptions in adjacent cells; test with a screen
	 reader such as NVDA or VoiceOver.}
\thesisaccessibilityrequirement{Colour and contrast}
	{Colour must not be the only carrier of meaning -- the Guide's worked
	 example supplements colour with patterns -- and contrast must be measured
	 at the intended font size.}
\thesisaccessibilityrequirement{Typeface}
	{A TrueType font recommended by ProQuest, at 10, 11 or 12 points.
	 Superscripts and subscripts no more than 2 points smaller than the body.}
